\section{Discussion}\label{sec:discussion}


\noindent \textbf{Limitation.} 
\toolname{} is designed to detect malicious behaviors that a model artifact executes against the host system by observing runtime interactions between the model process and the operating system. It cannot address attacks that target the model itself without affecting the host system, such as neural backdoor injection, where malicious functionality is embedded into learned parameters. These attacks compromise prediction integrity rather than runtime behavior, producing only normal tensor computations that generate no anomalous system calls or suspicious API invocations. Defending against such threats requires complementary techniques including model provenance verification, weight integrity checking, and robust training, which are orthogonal to and compatible with \toolname{}'s host-level defense. Dynamic detection also shares a general limitation. Conditionally triggered malicious logic runs only when its trigger condition is met during testing, so logic that is never activated goes undetected. Our threat model therefore recommends exercising each model across all anticipated usage scenarios during testing, including varied input shapes and edge-case inputs, to activate conditionally triggered logic and maximize coverage. We also plan a larger benign evaluation on public model hubs to stress the false-positive rate further, since the current results are based on the benign set described above.

\noindent \textbf{Relationship with static detection.} Static scanning tools and \toolname{} serve complementary roles in supply chain defense. Static scanners analyze model artifacts without executing them, making them lightweight and well-suited for platform-level screening of uploaded models. However, because they lack runtime context, they cannot reason about API usage semantics and are fundamentally blind to framework-native attacks that use only legitimate operations. \toolname{} operates at a deeper level by executing the model in a sandboxed environment to observe actual runtime behavior, enabling it to detect threats that evade static analysis. The two approaches can be combined in a defense-in-depth strategy: static scanning serves as an efficient first-pass filter to eliminate known exploit patterns, while \toolname{} provides thorough pre-deployment vetting for models that pass the initial screen.


\section{Conclusion}\label{sec:conclusion}

Malicious model artifacts present a growing threat to the AI supply chain: attackers embed harmful payloads within serialized models that execute during loading or inference, evading static scanners by using legitimate framework APIs. In this paper, we propose \toolname{}, a hybrid runtime defense system that combines non-bypassable kernel-level system-call monitoring with LLM-assisted Python semantic instrumentation to detect malicious behaviors in model artifacts. \toolname{} captures actual runtime execution evidence through a two-stage architecture: an unsupervised Autoencoder prefilter that surfaces anomalous low-level operations, and a cross-layer LLM reasoning engine that aligns these operations with their high-level semantic context to determine malicious intent. We evaluate \toolname{} on the MalHug Dataset and advanced synthetic attacks. The results show that \toolname{} achieves near-perfect detection with only 5.69\% average runtime overhead, demonstrating that it provides a practical, deployable defense for the AI model supply chain.